\section{Stage~H1 -- Service Identity and Mutual TLS Hardening}
\label{sec:rq21}

\subsection{Stage Objective}

As the communication-hardening stage of RQ2, this stage addresses the following question:

\begin{quote}
\textit{What performance and operational overhead is introduced when the selected
AKS microservice inference path is hardened with service identity, mutual TLS, and
inter-service authorization policy?}
\end{quote}

Stage~H1 adds the first security-hardening layer to the Azure chain-family deployment
carried forward from Stage~K and Stage~C, rather than reopening the local split-selection
result from Stages~L1--L2. The primary thesis-facing topology is
\texttt{chain\_2svc}, because it was the lowest-overhead non-monolithic chain in the
preceding Kubernetes and AKS chain-count stages.
The additional \texttt{chain\_5svc} run is used as a stress case to examine whether
the same security mechanism becomes more costly when the number of protected service
boundaries increases. The goal is not to evaluate trusted execution environments in
this stage, but to measure the cost of communication-level hardening: service identity,
mutual TLS, and service-account-based authorization for inter-service traffic.

The literature framing for service-mesh-based mTLS, identity, and
authorisation policy is discussed in \cref{sec:rw-security}.

\subsection{Experimental Setup}

The Stage~H1 benchmark used two frozen paired AKS artifact sets: a chain-2 main
driver \texttt{rq2\_1\_paired\_20260515\_160012} and a chain-5 maximum-depth
stress driver \texttt{rq2\_1\_paired\_20260517\_114303}. Both runs used the
same AKS cluster (\texttt{thesis-rq15}), strict same-node placement, and an
isolated system node pool so that control-plane pods did not share the
benchmark node.

All Stage~H1 conditions used the same pinned container image digest, shown across
two lines for typesetting:
\begin{center}
\small\texttt{sha256:d292aef407ee8a15da2bdaa680de1cc2}\\
\small\texttt{4b0975e7de81d27e1fc3a7d128b8736b}
\end{center}
The mTLS condition used the managed AKS Istio add-on
(\texttt{asm-1-29}). Each paired pass executed both security conditions; per
\texttt{paired\_summary.md} the alternation was mTLS-first on passes 1, 3, 5 and
plain-first on passes 2, 4 in both topologies, so each condition has five
executions and 1,000 measured iterations in total.

The complete benchmark configuration is reported in
\cref{tab:rq21_config} (\cref{sec:appendix-benchmark-configs}). The salient
parameters are the two paired topologies (\texttt{chain\_2svc} as the primary,
\texttt{chain\_5svc} as the maximum-depth stress case), the managed-Istio mesh
revision \texttt{asm-1-29}, a per-service budget of 1~CPU / 1~GiB plus an
additional per-service sidecar budget (100m--500m~CPU, 128--512~MiB) in the mTLS
condition, and five paired passes of 200 iterations per condition (1,000 after
merging) at seed~42.

The plain condition used Kubernetes service-to-service communication without mesh
injection. The mTLS condition added one Istio proxy per inference service, service
accounts for the protected services, workload-level peer-authentication objects, and
authorization-policy objects. The benchmark client remained outside the mesh.
Consequently, Stage~H1 measures hardening of the inter-service path inside the inference
chain rather than external ingress hardening. For \texttt{chain\_2svc}, the downstream
service was protected by strict mTLS and an authorization policy that allowed only the
expected upstream service account. For \texttt{chain\_5svc}, the same pattern was
applied to each downstream segment, producing four immediate-upstream authorization
policies.

Security validation passed in both frozen artifact sets. The mTLS runs passed static
manifest validation and runtime validation, including the expected
peer-authentication and authorization-policy counts. The full-chain
positive probe passed for every mTLS pass, and direct non-mesh probes were blocked for
the protected downstream segments. Resource metrics were complete in both runs.

\subsection{Results}

\begin{table}[ht]
\centering
\caption[Stage~H1 latency: plain vs.\ mTLS AKS chains]{Stage~H1 latency results for plain and mTLS AKS chains. Sources: \texttt{merged/aggregated\_results.md} in \texttt{rq2\_1\_paired\_20260515\_160012} (\texttt{chain\_2svc}) and \texttt{rq2\_1\_paired\_20260517\_114303} (\texttt{chain\_5svc}).}
\label{tab:rq21_latency}
\resizebox{\linewidth}{!}{%
\begin{tabular}{llccccc}
\toprule
\textbf{Topology} & \textbf{Condition} & \textbf{n} & \textbf{Mean (ms)} & \textbf{Median (ms)} & \textbf{p95 (ms)} & \textbf{Inferred non-compute/platform (ms)} \\
\midrule
\texttt{chain\_2svc} & \texttt{plain} & 1000 & 75.557 & 75.262 & 80.047 & 4.559 \\
\texttt{chain\_2svc} & \texttt{mtls}  & 1000 & 79.045 & 78.590 & 83.670 & 7.891 \\
\texttt{chain\_5svc} & \texttt{plain} & 1000 & 82.302 & 81.818 & 87.502 & 10.801 \\
\texttt{chain\_5svc} & \texttt{mtls}  & 1000 & 91.139 & 90.823 & 96.485 & 19.949 \\
\bottomrule
\end{tabular}
}
\end{table}

For the primary \texttt{chain\_2svc} topology, the mTLS/AuthZ hardening bundle
increased mean latency from 75.557~ms to 79.045~ms. The mean overhead was
therefore 3.488~ms, or 4.62\%, and p95 latency increased by 3.623~ms. The
inferred non-compute/platform component increased from 4.559~ms to 7.891~ms.
This supports the interpretation that most of the added latency is associated
with the communication/platform path introduced by the mesh rather than with a
change in neural-network computation.

For the additional \texttt{chain\_5svc} stress topology, the same hardening
bundle increased mean latency from 82.302~ms to 91.139~ms. The mean overhead
was 8.837~ms, or 10.74\%, and p95 latency increased by 8.983~ms. The inferred
non-compute/platform component increased from 10.801~ms to 19.949~ms. The
stress run therefore shows that the measured communication-hardening cost is
larger in the four-boundary stress case than in the selected two-service chain,
even though the main Stage~H1 answer remains anchored in \texttt{chain\_2svc}. This
two-point comparison should be read as depth sensitivity, not as a calibrated
scaling law.

\begin{table}[ht]
\centering
\caption[Stage~H1 paired-pass latency deltas]{Stage~H1 paired-pass latency deltas. Sources: \texttt{merged/aggregated\_results.md} (per-pass tables) in \texttt{rq2\_1\_paired\_20260515\_160012} and \texttt{rq2\_1\_paired\_20260517\_114303}. The pass-delta std is computed from the five per-pass deltas.}
\label{tab:rq21_paired_deltas}
\resizebox{\linewidth}{!}{%
\begin{tabular}{lccccc}
\toprule
\textbf{Topology} & \textbf{Mean delta (ms)} & \textbf{Mean delta (\%)} & \textbf{Min--max pass delta (ms)} & \textbf{Pass-delta std (ms)} & \textbf{p95 delta (ms)} \\
\midrule
\texttt{chain\_2svc} & 3.488 & 4.62  & 2.223--5.230 & 1.091 & 3.623 \\
\texttt{chain\_5svc} & 8.837 & 10.74 & 7.623--9.985 & 0.881 & 8.983 \\
\bottomrule
\end{tabular}
}
\end{table}

The paired-pass deltas in \cref{tab:rq21_paired_deltas} are consistent with the
overall means. In \texttt{chain\_2svc}, every paired pass shows a positive
mTLS/AuthZ hardening overhead, with pass-level deltas between 2.223~ms and
5.230~ms. In \texttt{chain\_5svc}, every paired pass is also positive, with
deltas between 7.623~ms and 9.985~ms. The alternating execution order reduces
the risk that the result is explained only by slow drift during the run.

\begin{figure}[htbp]
\centering
\resizebox{0.92\textwidth}{!}{%
\begin{tikzpicture}[
    font=\small,
    axis/.style={thick, draw=black!70, -{Latex[length=2mm]}},
    grid/.style={draw=black!12, thin},
    meanbar/.style={draw=blue!65!black, fill=blue!18, rounded corners=2pt, thick},
    tailbar/.style={draw=green!45!black, fill=green!18, rounded corners=2pt, thick},
    label/.style={font=\scriptsize, text=black!75, align=center},
    vallabel/.style={font=\scriptsize\bfseries, text=black!75},
    legendlabel/.style={font=\scriptsize, text=black!75},
    note/.style={font=\scriptsize, text=black!65, align=center}
]

% y scale: y = overhead milliseconds * 0.55
% chain_2svc mean delta: 3.488 ms -> 1.918
% chain_2svc p95 delta:  3.623 ms -> 1.993
% chain_5svc mean delta: 8.837 ms -> 4.860
% chain_5svc p95 delta:  8.983 ms -> 4.941

\foreach \ms/\y in {0/0,2/1.1,4/2.2,6/3.3,8/4.4,10/5.5} {
    \draw[grid] (0,\y) -- (8.8,\y);
    \node[anchor=east, font=\scriptsize, text=black!65] at (-0.18,\y) {\ms};
}

\draw[axis] (0,0) -- (9.15,0);
\draw[axis] (0,0) -- (0,5.9);

\node[rotate=90, font=\small, text=black!75] at (-0.95,2.95) {mTLS/AuthZ latency delta (ms)};
\node[font=\small, text=black!75] at (4.4,-1.45) {Topology and protected inter-service boundaries};

\draw[meanbar] (1.15,0) rectangle (1.95,1.918);
\draw[tailbar] (2.10,0) rectangle (2.90,1.993);
\draw[meanbar] (5.55,0) rectangle (6.35,4.860);
\draw[tailbar] (6.50,0) rectangle (7.30,4.941);

\node[vallabel, text=blue!65!black, above] at (1.55,1.918) {3.488};
\node[vallabel, text=green!45!black, above] at (2.50,1.993) {3.623};
\node[vallabel, text=blue!65!black, above] at (5.95,4.860) {8.837};
\node[vallabel, text=green!45!black, above] at (6.90,4.941) {8.983};

\node[label] at (2.025,-0.45) {\texttt{chain\_2svc}\\1 protected boundary\\+4.62\% mean};
\node[label] at (6.425,-0.45) {\texttt{chain\_5svc}\\4 protected boundaries\\+10.74\% mean};

\draw[meanbar] (0.65,5.18) rectangle (1.05,5.43);
\node[legendlabel, anchor=west] at (1.18,5.305) {Mean delta};

\draw[tailbar] (0.65,4.78) rectangle (1.05,5.03);
\node[legendlabel, anchor=west] at (1.18,4.905) {p95 delta};

\node[note, text width=8.6cm] at (4.4,-2.10) {
    Deltas are computed as mTLS minus plain within the paired AKS benchmark design.
};

\end{tikzpicture}

}
\caption[Stage~H1 mTLS/AuthZ latency overhead]{Stage~H1 mTLS/AuthZ latency overhead for the primary and stress AKS chains. Mean and p95 deltas are computed within the paired plain/mTLS design. The stress topology contains four protected inter-service boundaries, so the larger overhead visualises the observed stress-case sensitivity of sidecar-mediated communication hardening.}
\label{fig:rq21_mtls_latency_overhead}
\end{figure}

\Cref{fig:rq21_mtls_latency_overhead} visualises the paired latency deltas from
\cref{tab:rq21_paired_deltas}. The figure focuses on deltas rather than raw latency so
that the two Stage~H1 findings remain visible: the selected two-service chain pays a
modest overhead, while the deeper stress chain pays a larger overhead once four
inter-service boundaries are protected.

\begin{table}[ht]
\centering
\caption[Stage~H1 communication-hardening ablation]{Stage~H1 \texttt{chain\_2svc} ablation of the communication-hardening bundle. Source: \texttt{merged/rq2\_1\_ablation\_summary.md} in \texttt{rq2\_1\_ablation\_20260516\_112348}. As documented in the evidence manifest, ablation deltas are internally comparable only and must not be spliced into the frozen two-condition paired result.}
\label{tab:rq21_ablation}
\resizebox{\linewidth}{!}{%
\begin{tabular}{llccc}
\toprule
\textbf{Step} & \textbf{Condition} & \textbf{Mean latency (ms)} & \textbf{Adjacent delta (ms)} & \textbf{Interpretation} \\
\midrule
C0 & Plain Kubernetes, no mesh & 69.901 & -- & Non-mesh baseline \\
C1 & Mesh-default sidecars / auto-mTLS / no AuthZ & 72.470 & +2.569 & Main mesh data-path cost \\
C2 & Strict downstream mTLS / no AuthZ & 72.784 & +0.315 & Small added enforcement cost \\
C3 & Strict downstream mTLS plus AuthZ & 72.740 & -0.044 & Within pass-level noise \\
\bottomrule
\end{tabular}
}
\end{table}

An additional \texttt{chain\_2svc} ablation was run after the frozen paired
benchmark to interpret the hardening bundle more precisely. This ablation is
supporting evidence rather than a replacement for the frozen paired result; the
evidence manifest states that the ablation campaign is internally comparable
only and that its deltas must not be spliced into the frozen two-condition
chain-2 / chain-5 numbers. With that caveat, the full hardened ablation
condition reproduces the frozen result's order of magnitude: C3$-$C0 adds
2.840~ms, in the same range as the frozen paired \texttt{chain\_2svc} overhead
of 3.488~ms. As shown in \cref{tab:rq21_ablation}, most of the adjacent latency
change occurs when moving from plain Kubernetes to the mesh-default condition.
Strict mTLS enforcement over mesh-default traffic and the AuthorizationPolicy
step are both small relative to pass-level variation. The correct
interpretation is therefore that the deployed managed-Istio data path
introduces most of the measured overhead, while strict mTLS and identity-based
authorization mainly add enforcement semantics with little additional
steady-state latency in this workload.

\begin{table}[ht]
\centering
\caption[Stage~H1 resource and operational overheads]{Stage~H1 resource and operational overheads. Sources: \texttt{merged/paired\_summary.md} in \texttt{rq2\_1\_paired\_20260515\_160012} (\texttt{chain\_2svc}) and \texttt{rq2\_1\_paired\_20260517\_114303} (\texttt{chain\_5svc}).}
\label{tab:rq21_resource_overhead}
\resizebox{\linewidth}{!}{%
\begin{tabular}{lcccccc}
\toprule
\textbf{Topology} & \textbf{Sidecar CPU (mCPU)} & \textbf{Total pod CPU delta (mCPU)} & \textbf{Sidecar memory (MiB)} & \textbf{Total pod memory delta (MiB)} & \textbf{Schedule-to-ready delta (s)} & \textbf{Additional objects} \\
\midrule
\texttt{chain\_2svc} & 34.684 & $-$56.813 & 82.125  &  93.819 & 6.30 & 5 \\
\texttt{chain\_5svc} & 87.602 &   37.490  & 223.714 & 264.143 & 6.84 & 14 \\
\bottomrule
\end{tabular}
}
\end{table}

The mTLS condition also increased resource requests and deployment complexity.
For \texttt{chain\_2svc}, the mesh added two injected proxies, two service
accounts, two peer-authentication objects, one authorization-policy object,
200m requested sidecar CPU, and 256~MiB requested sidecar memory. For
\texttt{chain\_5svc}, the mesh added five injected proxies, five service
accounts, five peer-authentication objects, four authorization-policy objects,
500m requested sidecar CPU, and 640~MiB requested sidecar memory. In both
cases, the always-on Istio control-plane pods were present on the isolated
system pool rather than the benchmark node. The CPU counter deltas should be
read as window-level operational measurements rather than perfect causal
attribution: in \texttt{chain\_2svc}, measured application-container CPU was
lower in the mTLS condition, which makes the total pod CPU delta negative
despite a clearly positive sidecar contribution.

The interpretation of these results and their contribution to RQ2 are presented in
\cref{sec:analysis-rq21}.

\section{Stage~H1 (Split Sensitivity) -- Single-Hop mTLS vs Activation Size}
\label{sec:rq21-splitsens}

This stage corresponds to the methodological design declared in
\cref{sec:methods-rq21-splitsens}: a supporting campaign that varies the
single protected hop's activation size across \texttt{layer1},
\texttt{layer2}, \texttt{layer3}, and \texttt{layer4} and pairs each split
point's plain and mTLS conditions, in order to measure whether the
single-hop mTLS overhead scales with the activation crossing that hop.

\begin{quote}
\textbf{TODO (author).} The split-sensitivity campaign results
(per-split mean and p95 mTLS overheads, per-split operational overhead, and
per-split resource summary derived from
\texttt{rq2\_1\_mtls\_split\_20260516\_095406}) are not yet written into this
chapter. The appendix entry for this campaign is also marked as pending in
\cref{ch:appendix} (see the planning comment at the top of
\texttt{appendix/appendix.tex}). Until the per-split numbers are inserted
here -- or moved into an appendix subsection and forward-referenced from
this stub -- the split-sensitivity finding declared in
\cref{sec:methods-rq21-splitsens} has no visible Results entry in this
chapter and must not be cited as a measured outcome. This stub is
intentionally a placeholder so that the gap is visible at the
table-of-contents level and is not silently absorbed into the primary Stage~H1
result.
\end{quote}

\section{Stage~H1 (Multi-Node Alternative) -- Inter-Node Sensitivity of mTLS Overhead}
\label{sec:rq21b}

This stage corresponds to the methodological design declared in
\cref{sec:methods-rq21b}: a paired sensitivity stage that forces the two
\texttt{chain\_2svc} inference services onto distinct AKS nodes and re-pairs
the plain and mTLS conditions under the same managed-Istio configuration as
the primary Stage~H1 run. Every paired execution passed the multi-node placement
contract described in \cref{sec:methods-rq21b}.

Under the multi-node placement, the mean mTLS overhead on \texttt{chain\_2svc}
is +2.71~ms, or +3.90\%, with every paired execution passing multi-node
placement validation. This is in the same direction and order of magnitude as
the +3.488~ms / +4.62\% single-node primary result reported in
\cref{tab:rq21_latency}, and is slightly smaller in percent terms, consistent
with a larger plain baseline absorbing more of the proxy cost once inter-node
transport is added. The result is treated as a sensitivity envelope around the
primary single-node Stage~H1 result, not as a substitute for it.

The interpretation of this multi-node sensitivity result, and its role in
bounding the single-node primary number, are developed under
\cref{sec:analysis-rq21} as Finding~7.

\section{Stage~H2 -- Confidential VM Execution Hardening}
\label{sec:rq22}

\subsection{Stage Objective}

As the confidential-execution stage of RQ2, this stage addresses the following question:

\begin{quote}
\textit{What additional performance and operational overhead is introduced when
VM-level confidential execution is added to the already communication-secured
AKS inference chain, when the downstream protected service is placed on AMD
SEV-SNP confidential nodes?}
\end{quote}

Stage~H2 adds a second security-hardening layer to the Azure deployment path. Stage~H1
already established the communication-secured baseline: the selected
\texttt{chain\_2svc} inference pipeline runs on AKS with managed-Istio mTLS,
service identity, and an authorisation policy on the protected downstream
service. Stage~H2 keeps that communication-security contract fixed and changes the
execution environment of the downstream service. The confidential condition
places \texttt{service2} on an Azure confidential VM node pool backed by AMD
SEV-SNP, while the standard condition uses matched non-confidential AMD VM
nodes for both services. Per the canonical evidence manifest, the
thesis-facing Stage~H2 scope is \texttt{service2\_only}; placing both services on
confidential nodes (full two-service TEE) is treated as future work and is not
evaluated here. He~et~al.~\cite{he2021_attacking_collaborative} motivate
protecting the activations received by the downstream segment beyond the
in-transit mTLS layer, although they evaluate a different (malicious
remote-service) threat model than the SEV-SNP infrastructure-level protection
measured here.

The purpose of this stage is therefore incremental. It does not ask whether
mTLS is useful; that was answered in Stage~H1. Instead, it asks what extra cost is
paid when runtime protection against an additional infrastructure-level threat model is
added on top of the already secured service-to-service path.

The literature framing for AMD SEV-SNP, VM-level confidential computing, and
its interaction with sidecar-based service meshes is discussed in
\cref{sec:rw-security}.

\subsection{Experimental Setup}

Stage~H2 used the frozen rolling AKS artifact \texttt{rq2\_2\_confidential\_20260516\_135726}.
The run preserves the Stage~H1 managed-Istio mTLS and authorisation semantics, and
collects newly paired standard baselines on \texttt{Standard\_D8as\_v5} nodes
rather than reusing the Stage~H1 numbers. The application image is the same ACR
digest used in Stage~C and Stage~H1
(\texttt{thesisrq15acr.azurecr.io/thesis-inference\\@sha256:d292aef407ee\dots}),
recorded in \texttt{image\_reference.txt}.

The complete benchmark configuration is reported in
\cref{tab:rq22_config} (\cref{sec:appendix-benchmark-configs}). The salient
parameters are the standard VM size \texttt{Standard\_D8as\_v5} and the
confidential \texttt{Standard\_DC8as\_v5} (AMD SEV-SNP), a \texttt{service2\_only}
TEE scope, the \texttt{chain\_2svc} topology inheriting the Stage~H1 mTLS and
authorisation semantics, and five paired passes of 200 iterations per condition
(1,000 after merging).

The benchmark compared two conditions. In the standard condition, service1 and
service2 both ran on \texttt{Standard\_D8as\_v5}. In the confidential
condition, service1 remained on \texttt{Standard\_D8as\_v5}, while service2 ran
on \texttt{Standard\_DC8as\_v5} backed by AMD SEV-SNP. This isolates the cost
of protecting the downstream service that receives intermediate activations
and executes the later ResNet-18 segment. The rolling design alternates
standard and confidential passes from a seeded order (standard-first on passes
1, 3, 5; confidential-first on passes 2, 4 per
\texttt{execution\_plan.json}), and parity validation against the monolithic
reference (\texttt{max\_abs\_diff}~$=0.0$, \texttt{num\_inputs}~$=5$) passed in
every per-pass \texttt{parity\_validation.json}.

\subsection{Results}

\begin{table}[ht]
\centering
\caption[Stage~H2 service2-only confidential execution results]{Stage~H2 service2-only confidential execution results. Source: per-pass \texttt{benchmark/raw\_iterations.csv} in \texttt{rq2\_2\_confidential\_20260516\_135726/conditions/}, aggregated across the five standard and five confidential passes (1,000 iterations per condition).}
\label{tab:rq22_service2_results}
\resizebox{\linewidth}{!}{%
\begin{tabular}{lccccc}
\toprule
\textbf{Condition} & \textbf{Service2 VM} & \textbf{n} & \textbf{Mean (ms)} & \textbf{Median (ms)} & \textbf{p95 (ms)} \\
\midrule
\texttt{standard}     & \texttt{Standard\_D8as\_v5}  & 1000 & 58.492 & 58.289 & 60.255 \\
\texttt{confidential} & \texttt{Standard\_DC8as\_v5} & 1000 & 60.346 & 59.756 & 63.856 \\
\midrule
\textbf{Delta} & --- & --- & \textbf{+1.854 / +3.17\%} & \textbf{+1.467 / +2.52\%} & \textbf{+3.601 / +5.98\%} \\
\bottomrule
\end{tabular}
}
\end{table}

In the service2-only TEE run, mean end-to-end latency increased from 58.492~ms
to 60.346~ms. The mean overhead was therefore 1.854~ms, or 3.17\%. Median
latency increased by 1.467~ms, and p95 latency increased by 3.601~ms.
Sequential throughput, computed from the mean latency, decreased from
17.096~req/s to 16.571~req/s, a reduction of approximately 3.07\%.

This is the only confidential-execution scope evaluated in the thesis. The
evidence manifest scopes Stage~H2 to \texttt{service2\_only} and explicitly treats
a full two-service confidential placement as future work. The trade-off table
in \cref{sec:rq23} therefore compares plain AKS, mTLS, and mTLS plus selective
TEE only.

The interpretation of these results and their contribution to RQ2 are presented in
\cref{sec:analysis-rq22}.

\section{Trade-off Synthesis -- Security-Hardening Recommendation}
\label{sec:rq23}

\subsection{Stage Objective}

As the concluding trade-off synthesis of RQ2, this stage addresses the following question:

\begin{quote}
Which security-hardening configuration provides the most appropriate trade-off
between performance cost, operational complexity, and protection scope for the
selected AKS microservice inference deployment?
\end{quote}

The trade-off synthesis is a synthesis stage rather than a new benchmark. Stage~H1 measured the
cost of adding service identity, mutual TLS, and inter-service authorization
policy to the selected Azure inference path. Stage~H2 then measured the additional
cost of placing the downstream protected service on an AMD SEV-SNP confidential
VM node pool while keeping the communication-security contract fixed. The trade-off synthesis
combines these results to compare the evaluated hardening levels and identify
when each level is appropriate.

\subsection{Synthesis Basis}

The synthesis compares three deployment configurations evaluated in this
thesis. The first is the plain AKS baseline from Stage~C, where the inference
chain runs without service-mesh mTLS or confidential-node placement. The second
is the Stage~H1 communication-secured configuration, where managed-Istio mTLS,
service identity, and authorization policy are added to the selected chain.
The third is the Stage~H2 selective confidential-execution configuration, where
only the downstream protected service runs on an AMD SEV-SNP confidential
node. Extending the confidential-computing boundary to both services (full
two-service TEE) is out of the thesis-facing scope per the evidence manifest
and is discussed only as future work.

These configurations should not be interpreted as interchangeable variants of
the same security mechanism. They protect different parts of the system. Plain
AKS provides the performance baseline but does not add the communication-level
identity and encryption policy evaluated in Stage~H1. The mTLS/AuthZ configuration
hardens the service-to-service path by authenticating workloads and enforcing
authorization policy. The confidential-node configuration adds a VM-level
runtime protection boundary for the downstream service placed on AMD SEV-SNP
infrastructure. As discussed in
\cref{sec:methods-rq22,sec:analysis-rq22,sec:analysis-threats}, the SEV-SNP
threat-model scope is bounded: residual attack surfaces such as malicious-host
interrupt injection remain outside the memory-encryption
guarantee~\cite{schluter2024_heckler}, and this thesis does not experimentally
evaluate such attacks.

\subsection{Trade-off Summary}

\begin{table}[t]
\centering
\caption[Trade-off synthesis of security-hardening levels]{Trade-off synthesis of evaluated security-hardening levels. mTLS overheads are taken from \cref{tab:rq21_latency}; the selective TEE overhead is taken from \cref{tab:rq22_service2_results} relative to its own paired standard-node mTLS baseline. Each cost is reported as both the mean delta and the p95 delta, because the two security mechanisms differ in how the latency cost is distributed between the centre and the tail of the distribution.}
\label{tab:rq23_security_tradeoff}
\begin{tabular}{>{\raggedright\arraybackslash}p{0.14\textwidth}>{\raggedright\arraybackslash}p{0.30\textwidth}>{\raggedright\arraybackslash}p{0.21\textwidth}>{\raggedright\arraybackslash}p{0.21\textwidth}}
\toprule
Configuration & Protection scope & Mean cost & p95 cost \\
\midrule
Plain AKS
& Baseline AKS deployment without mesh mTLS/AuthZ or confidential-node placement.
& Baseline (reference).
& Baseline (reference). \\
\addlinespace
mTLS/AuthZ
& Service identity, encrypted service-to-service communication, and authorization policy for the selected inference path.
& chain\_2svc: +3.488~ms / +4.62\%\newline chain\_5svc: +8.837~ms / +10.74\%.
& chain\_2svc: +3.623~ms / +4.53\%\newline chain\_5svc: +8.983~ms / +10.27\%. \\
\addlinespace
mTLS/AuthZ + selective TEE
& Communication hardening plus AMD SEV-SNP VM-level confidential execution for the downstream protected service (\texttt{service2\_only}).
& +1.854~ms / +3.17\% over the newly paired standard-node mTLS baseline.
& +3.601~ms / +5.98\% over the same baseline. \\
\bottomrule
\end{tabular}
\end{table}

The trade-off is not a single global ranking because the preferred configuration
depends on the threat model. If performance is the only criterion, the plain
AKS deployment is the best option. However, it does not provide the additional
service identity, encrypted inter-service communication, or
authorization-policy enforcement evaluated in Stage~H1. If the main security
requirement is to protect communication between inference services, the
mTLS/AuthZ configuration provides the most appropriate default balance. It
adds a mean latency overhead of 3.488~ms, or 4.62\%, for the selected
\texttt{chain\_2svc} topology, while enforcing the intended
service-account-based policy. The \texttt{chain\_5svc} stress case shows a larger
cost when four inter-service boundaries are protected, reaching 8.837~ms, or
10.74\%, for the full five-service chain. The mean and p95 overheads fall in the
same range, so the mTLS cost is centred rather than concentrated in the latency
tail.

If the threat model includes host-level exposure of the downstream inference
component, mTLS alone is not sufficient because it protects the communication
path rather than the runtime memory of the service, and prior work shows that
an untrusted remote service holding only intermediate feature maps can in
principle recover the input~\cite{he2021_attacking_collaborative}. Selective
confidential execution provides a narrower hardening step against an
infrastructure-level adversary on the host platform. The service2-only
confidential-node run increases mean latency by 1.854~ms, or 3.17\%, relative
to its paired standard-node mTLS baseline. In contrast to mTLS, this cost is
tail-weighted: the p95 delta of 3.601~ms (5.98\%) is roughly 1.9~times the mean
delta in percentage terms, so tail-sensitive service-level objectives should
budget against the p95 figure. This makes selective TEE attractive
when the downstream service is the main infrastructure-level sensitivity
boundary, while bounded by SEV-SNP's known residual attack
surfaces~\cite{schluter2024_heckler}, which this thesis does not
experimentally evaluate.

Extending the confidential-execution boundary to both services would broaden
the runtime-protection scope at additional cost. That configuration is out of
the thesis-facing scope of Stage~H2 (see the evidence manifest) and is therefore
discussed as future work rather than as a measured trade-off point.

The interpretation of these results and their contribution to RQ2 are presented in
\cref{sec:analysis-rq23}.
